\documentclass{article}%
\usepackage[T1]{fontenc}%
\usepackage[utf8]{inputenc}%
\usepackage{lmodern}%
\usepackage{textcomp}%
\usepackage{lastpage}%
\usepackage{geometry}%
\geometry{margin=2.2cm}%
%
\usepackage{amsmath}%
\usepackage{amssymb}%
\usepackage{booktabs}%
\usepackage{longtable}%
\usepackage{array}%
\usepackage{hyperref}%
\hypersetup{hidelinks}%
\title{ITA 2026 {-} Física (claude\_haiku\_5\_5) x Poliedro\\[4pt]\large Comparação com o gabarito Poliedro Resolve\thanks{Gerado com gpt-resolve: \url{https://github.com/lgabs/gpt-resolve}}}%
\date{10/11/2026}%
%
\begin{document}%
\normalsize%
\maketitle%
\section*{Resumo}\textbf{Nota média: 9,5/10} (21 de 22 itens com a resolução correta; cada item vale o mesmo dentro da questão).\par\medskip%
\begin{longtable}{@{}c c p{0.78\textwidth}@{}}\toprule Questão & Nota & Resumo \\ \midrule
1 & 10,0 & Os três itens concordam com o gabarito do Poliedro Resolve: dois em forma equivalente e um idêntico. No item (a), v$^2$(1 - $\mu$ tan$\theta$)/{[}g($\mu$ + tan$\theta$){]} coincide com (v$^2$/g)(cos$\theta$ - $\mu$ sen$\theta$)/(sen$\theta$ + $\mu$ cos$\theta$). No item (b), v$^2$/(g tan$\theta$) é o mesmo que (v$^2$/g)$\cdot$cotg $\theta$. No item (c), as duas dão v$^2$/($\mu$g). Refiz as contas e conferi a equivalência numericamente. Também confirmei, ampliando a figura, que a face de maior diâmetro das rodas fica voltada para dentro, o que justifica a orientação das normais adotada por ambos. A resolução está correta e é mais completa que o gabarito: prova que a configuração limite é de fato o máximo e discute o caso $\mu$ tan$\theta$ $\geq$ 1, que o gabarito omite. Não há erros a apontar. \\
2 & 10,0 & A resolução e o gabarito do Poliedro concordam na Questão 2. No item a, as normais são idênticas; conferi as fórmulas e o exemplo numérico. No item b, os dois obtêm f\_A + f\_B = (M + m + m\_c)$\cdot$g$\cdot$sen$\theta$ e reconhecem que a repartição entre os apoios é estaticamente indeterminada. A resolução acrescenta valores individuais (f\_i = N\_i$\cdot$tan$\theta$) sob uma hipótese que ela declara, sem afirmar que eles decorrem da estática; por isso o item b foi classificado como equivalente. O que fica para um professor é saber o que a banca esperava no item b, já que o enunciado pede o atrito "em cada apoio" e o problema é indeterminado. Não encontrei erro de conta nem de física em nenhum dos dois. \\
3 & 10,0 & A resolução está certa nos dois itens e bate com o gabarito do Poliedro. Item (a): N\_r = Mg cos$\theta$ + M$\omega$$^2$(L - 2a)/2, ao longo da haste, apontando para o trecho de comprimento a. Item (b): $\theta$ = arccos{[}$\omega$$^2$(2a - L)/(2g){]}. Conferi numericamente: o teorema do centro de massa no referencial inercial (usado na resolução) e as forças centrífugas no referencial girante (usadas no gabarito) dão a mesma expressão. A resolução vai além do gabarito: dá a condição de existência $\omega$$^2$|L - 2a| $\leq$ 2g, trata o caso a $\geq$ L/2 e acerta a componente tangencial na observação, F\_t = MgL$^2$ sin$\theta$/{[}L$^2$ + 3(L - 2a)$^2${]}. Não há divergência que mude a nota. \\
4 & 10,0 & A resolução e o gabarito do Poliedro concordam nos dois itens. No item a, os dois dão y$_{2}$ = 2b ou y$_{2}$ = -4b. No item b, a resposta da resolução, -(GMm/5){[}16/$\surd$(25a$^2$ - 15b$^2$ - 24ac) + 1/$\surd$(2a$^2$ + 3b$^2$ - 2ac){]}, é a mesma expressão do gabarito, -GMm$_{2}${[}4/d$_{1}$ + 1/d$_{2}${]}, só que expandida: m$_{2}$ = m/5 e c$^2$ = a$^2$ - b$^2$. Conferi isso algebricamente e numericamente; as duas escolhem a estrela no foco (c,0) e y$_{2}$ = 2b. A resolução está correta dentro do modelo pedido, e também estão certos a tabela de força bruta, o caso circular e as demais verificações. O ponto que pode mudar a nota é o próprio enunciado: a premissa de que o CM continua na órbita original é fisicamente falsa para uma separação da ordem de a. O Poliedro sugere anular a questão, e a resolução também aponta a falha. \\
5 & 5,0 & No item a, a resolução e o gabarito dão o mesmo resultado: W = 70(M1 - M2)g$\alpha$L = 700$\alpha$L(M1 - M2) J. No item b, o método é o mesmo ($\eta$ = W/Q\_q, com água, cobre por Dulong{-}Petit e trabalho sobre M1). A diferença é que a expressão final da resolução não tem a parcela m$\cdot$g$\cdot$$\alpha$$\cdot$L/2 da energia potencial da própria viga, que o gabarito inclui corretamente. A resolução discute esse termo e o descarta como desprezível (\textasciitilde{}10$^{-}$$^{7}$ do termo do cobre; $\eta$ numericamente igual), mas a omissão é incoerente porque M1g$\alpha$L, da mesma ordem, foi mantido. Veredito: item a correto; item b certo na física e quase completo, com o gabarito mais completo e possível desconto pequeno. \\
6 & 10,0 & A resolução está certa nos três itens e concorda com o gabarito Poliedro Resolve. Conferi o enunciado no PDF (página 7), e a transcrição está fiel.\newline%
{-} a) Os dois dão L = 85 cm.\newline%
{-} b) Os dois dão G = 16 em módulo (-16 com sinal) e imagem invertida.\newline%
{-} c) A objetiva forma uma imagem virtual a 80 cm (p1' = -80 cm), que serve de objeto para a ocular a p2 = 165 cm. A imagem final é real e invertida e fica a 165/32 $\approx$ 5,16 cm da ocular, com m = -1/16. É menor que o objeto. O gabarito chega ao mesmo resultado, só que arredondado: 5,2 cm e -0,063.\newline%
Refiz todas as contas com frações exatas e encontrei os mesmos valores. Não há erro nem no gabarito nem na resolução. \\
7 & 10,0 & Questão 7: a resolução bate com o gabarito do Poliedro nos dois itens. No item a, os dois dão I = x\^{}\{N+1\} I\_0 cos\^{}\{2N\}($\varphi$/N). No item b, dão r = $\surd$(P/($\pi$ x\^{}\{N+1\} I\_0 cos\^{}\{2N\}($\varphi$/N))). O enunciado transcrito confere com a página 8 do PDF. A física está correta: lei de Malus com $\beta$ = $\varphi$/N entre polarizadores consecutivos, nenhum fator de cosseno no primeiro polarizador e perda x em cada um dos N+1 polarizadores. Conferi em Python todos os exemplos numéricos da resolução e estão certos. Não há divergência. As dúvidas listadas são só de interpretação (como aplicar x, o papel do analisador) e de completude. \\
8 & 10,0 & A resolução e o gabarito coincidem nos dois itens: a) P\_min = 35 W; b) $\Delta$t\_max = (56/25)$\pi$$^2$$\times$10$^{-}$$^{6}$ s $\approx$ 2,2$\times$10$^{-}$$^{5}$ s, porque 2,24$\pi$$^2$$\times$10$^{-}$$^{6}$ = (56/25)$\pi$$^2$$\times$10$^{-}$$^{6}$. Refiz todas as contas em python e chego aos mesmos números: E\_tot = 42 kJ, L = 8$\pi$$^2$$\times$10$^{-}$$^{7}$ H, $\Delta$t = 14L/5 $\approx$ 22,1 µs. A resolução vai além do gabarito ao mostrar que a cota energética dá o mesmo valor e ao justificar o máximo pelo argumento |M| $\leq$ L. O enunciado transcrito confere com o original (pág. 9). Há só pontos menores: o erro de unidade de $\mu$0 na prova, uma possível leitura alternativa de 'período de tempo $\Delta$t' e a hipótese H3, que não é necessária para chegar à resposta. \\
9 & 10,0 & A questão 9 está certa nos dois itens, e a resolução bate com o gabarito do Poliedro. Conferi o enunciado na página 10 do PDF. a) d = 600$\surd$2 nm $\approx$ 848,5 nm: a resolução escreve 6$\surd$2$\times$10$^{-}$$^{7}$ m, que é o mesmo valor em outra unidade. b) O violeta cai a 20$\surd$14/7 cm $\approx$ 10,7 cm do ponto central, com sen $\theta$' = $\surd$2/3. Os dois usam o mesmo raciocínio: ângulo pela tangente e sen $\theta$ na equação da grade. Refiz as contas em Python e os valores conferem. Não achei erro em nenhum dos dois, só duas observações menores sobre a leitura do enunciado e hipóteses implícitas. \\
10 & 10,0 & A resolução e o gabarito Poliedro estão de acordo nos dois itens e ambos estão certos. a) 387/32 = 12,09375 anos ($\approx$ 12,1 anos). b) 387/40 = 9,675 anos ($\approx$ 9,7 anos), tempo próprio da nave 2 com $\gamma$2 = 5/4. Refiz as contas com frações exatas em python. Fiz também uma checagem independente pela transformação de Lorentz no referencial da nave 2, que confirma os 9,675 anos. Não há erro em nenhum dos dois. Só fica uma ambiguidade pequena de interpretação ('tempo na Terra' como coordenada de tempo, e não como instante de observação, que daria 19,35 anos), e a resolução trata dela explicitamente. \\
\bottomrule\end{longtable}%
\section*{Principais pontos de dúvida}%
\begin{itemize}%
\item%
\textbf{Q2} \textit{(statement)} O item b pede o atrito "em cada apoio", mas, com dois apoios com atrito, o problema é estaticamente indeterminado. Não se sabe o que a banca do ITA esperava: só a soma (como o Poliedro), f\_i = N\_i$\cdot$tan$\theta$ (atrito proporcional à normal, como no fechamento da resolução) ou f\_A = f\_B. Isso pode mudar o critério de correção do item. A resolução cobre os dois primeiros casos, porque dá a soma e também a repartição sob uma hipótese declarada.%
\item%
\textbf{Q4} \textit{(statement)} A questão depende de uma premissa que o enunciado não declara: depois da cisão, o CM de m$_{1}$ + m$_{2}$ continuaria na elipse original. Isso só vale logo após a separação. Com os dados, a distância P$_{1}$P$_{2}$ = $\surd$(a$^2$ + 49b$^2$)/4 é comparável a a, e aí os fragmentos sentem campos bem diferentes. O Poliedro sugere anular a questão; a resolução também aponta a falha, mas segue o modelo. Convém confirmar se a banca anulou ou alterou a questão, porque isso muda a nota dos dois itens.%
\item%
\textbf{Q5} \textit{(solution)} No item b, a expressão final de $\eta$ não tem a parcela m$\cdot$g$\cdot$$\alpha$$\cdot$L/2 (energia potencial da viga, cujo centro de massa sobe $\Delta$L/2), que o gabarito inclui em Q\_q. A resolução reconhece o termo e o descarta como desprezível, mas mantém M1$\cdot$g$\cdot$$\alpha$$\cdot$L, que é da mesma ordem. Um corretor pode descontar parte da nota do item b por resposta literal incompleta.%
\end{itemize}%
\section*{Questão 1}%
\href{https://poliedroresolve.sistemapoliedro.com.br/vestibulares/ita/2026/ita-2-fase-dia-3-30-10-2025-2-fase-dia-3/dissertativa-1-fisica-2-fase-dia-3-ita-2026}{Resolução do gabarito}\par\smallskip%
\begin{longtable}{@{}p{0.06\textwidth} p{0.33\textwidth} p{0.33\textwidth} p{0.1\textwidth} p{0.08\textwidth}@{}}\toprule Item & Resolução & Gabarito & Veredito & Certo \\ \midrule
a & R\_min = v$^2$/(g$\cdot$tan($\theta$ + arctan $\mu$)) = v$^2$(1 - $\mu$ tan$\theta$)/{[}g($\mu$ + tan$\theta$){]}, para $\mu$ tan$\theta$ < 1 (para $\mu$ tan$\theta$ $\geq$ 1 o modelo não impõe cota inferior) & r\_min = (v$^2$/g)$\cdot$(cos$\theta$ - $\mu$ sen$\theta$)/(sen$\theta$ + $\mu$ cos$\theta$), no limite N\_i $\approx$ 0, Fat\_i $\approx$ 0, Fat\_E = $\mu$N\_E & equivalente & ambos \\
b & R\_min = v$^2$/(g tan$\theta$) & r\_min = (v$^2$/g)$\cdot$cotg $\theta$ & equivalente & ambos \\
c & R\_min = v$^2$/($\mu$g) & r\_min = v$^2$/($\mu$g) & igual & ambos \\
\bottomrule\end{longtable}%
\begin{itemize}%
\item%
\textbf{a)} Multiplicando o numerador e o denominador da forma da resolução por cos$\theta$, obtém{-}se exatamente a expressão do gabarito. Também vale (cos$\theta$ - $\mu$ sen$\theta$)/(sen$\theta$ + $\mu$ cos$\theta$) = 1/tan($\theta$+$\varphi$), com $\varphi$ = arctan $\mu$. Conferi numericamente para valores aleatórios de $\theta$ e $\mu$, e as três formas coincidem. A configuração limite é a mesma nas duas: N\_I = 0, f\_I = 0 e f\_E = $\mu$N\_E, com o atrito na roda externa voltado para o centro da curva e para baixo, o que é coerente com as equações (I) e (II) do gabarito. A resolução ainda prova que esse é o máximo de $\Sigma$H (o ângulo de cada reação com a vertical é no máximo $\theta$+$\varphi$) e trata à parte o caso $\mu$ tan$\theta$ $\geq$ 1, que o gabarito não discute. Nesse caso a fórmula do gabarito daria r $\leq$ 0.%
\item%
\textbf{b)} Como cotg $\theta$ = 1/tan$\theta$, as respostas são idênticas. Refiz a conta: (N\_E - N\_I) sen$\theta$ = Mv$^2$/R e (N\_E + N\_I) cos$\theta$ = Mg. Impondo N\_I = 0, sai v$^2$/(gR) = tan$\theta$. A figura, conferida no PDF ampliado, mostra a face de maior diâmetro voltada para dentro. Portanto a normal na roda externa aponta para o centro da curva, como as duas resoluções supõem.%
\item%
\textbf{c)} Com $\theta$ = 0 a normal é vertical e o atrito é horizontal, logo Mv$^2$/R $\leq$ $\mu$$\Sigma$N = $\mu$Mg e R $\geq$ v$^2$/($\mu$g). O gabarito chega ao mesmo valor fazendo $\theta$ = 0 em (a). A resolução observa que, com rodas cilíndricas, o limite vem do escorregamento lateral e não da perda de contato por N = 0. É uma nuance correta, que não muda a resposta.%
\end{itemize}%
\paragraph{Pontos de dúvida}%
\begin{itemize}%
\item%
\textit{(key)} A fórmula do item (a) do gabarito só vale para $\mu$ tan$\theta$ < 1, isto é, $\theta$ + arctan $\mu$ < 90$^\circ$. Para $\mu$ tan$\theta$ $\geq$ 1 ela dá r $\leq$ 0, que não tem sentido físico, e o gabarito não menciona essa restrição. A resolução trata esse caso (efeito de cunha, sem cota inferior no modelo rígido de Coulomb). Para rodas reais ($\theta$ pequeno) a restrição é irrelevante.%
\item%
\textit{(statement)} Gabarito e resolução tratam o trem como partícula: impõem apenas o equilíbrio de forças e ignoram o equilíbrio de momentos (tombamento), já que a altura do centro de massa não é dada. Na configuração limite com N\_I = 0, a reação da roda externa precisaria passar pelo CM. Essa é uma hipótese implícita do enunciado, adotada por ambos.%
\item%
\textit{(statement)} Com atrito nas duas rodas, a repartição das forças é estaticamente indeterminada. O critério 'raio mínimo' é interpretado, em ambas, como a máxima força centrípeta que os contatos conseguem fornecer (roda interna com N = 0 e roda externa no limite de atrito). É a interpretação padrão e esperada.%
\item%
\textit{(key)} A transcrição do diagrama do gabarito descreve Fat\_E 'para baixo/para fora'. Pelas equações do próprio gabarito (N\_E sen$\theta$ + Fat\_E cos$\theta$ = Mv$^2$/r e N\_E cos$\theta$ - Fat\_E sen$\theta$ = Mg), porém, a componente horizontal do atrito na roda externa aponta para o centro da curva, como na resolução. Provavelmente é só imprecisão da transcrição da imagem, sem efeito no resultado.%
\end{itemize}%
\section*{Questão 2}%
\href{https://poliedroresolve.sistemapoliedro.com.br/vestibulares/ita/2026/ita-2-fase-dia-3-30-10-2025-2-fase-dia-3/dissertativa-2-fisica-2-fase-dia-3-ita-2026}{Resolução do gabarito}\par\smallskip%
\begin{longtable}{@{}p{0.06\textwidth} p{0.33\textwidth} p{0.33\textwidth} p{0.1\textwidth} p{0.08\textwidth}@{}}\toprule Item & Resolução & Gabarito & Veredito & Certo \\ \midrule
a & N\_A = g cos$\theta$ {[}M(L/2 - d) + m(L - a - d) + m\_c(b - d){]}/(L - 2d); N\_B = g cos$\theta$ {[}M(L/2 - d) + m(a - d) + m\_c(L - b - d){]}/(L - 2d), com N\_A + N\_B = (M + m + m\_c) g cos$\theta$ & N\_A = {[}m(L - a - d) + M(L/2 - d) + m\_C(b - d){]}$\cdot$g$\cdot$cos$\theta$/(L - 2d); N\_B = {[}m(a - d) + M(L/2 - d) + m\_C(L - b - d){]}$\cdot$g$\cdot$cos$\theta$/(L - 2d) & igual & ambos \\
b & Pela estática, f\_A + f\_B = (M + m + m\_c) g sen$\theta$, e só a soma fica determinada. Sob a hipótese explícita de mobilização uniforme do atrito (f\_i/N\_i igual nos dois apoios, o que dá tan$\theta$), f\_A = N\_A tan$\theta$ e f\_B = N\_B tan$\theta$. A resolução diz expressamente que essa repartição não decorre da estática. & Só a soma é determinável: Fat\_A + Fat\_B = (m + M + m\_C) g sen$\theta$, para cima do plano. Faltam dados para obter cada atrito separadamente; o problema é indeterminado nesse ponto. & equivalente & ambos \\
\bottomrule\end{longtable}%
\begin{itemize}%
\item%
\textbf{a)} As expressões são idênticas, só muda a ordem dos termos. Conferi numericamente com valores aleatórios: a resultante normal e o momento em torno de um ponto qualquer da barra zeram. O exemplo da resolução também confere: N\_A $\approx$ 57,92 N, N\_B $\approx$ 54,67 N e soma 112,58 N = 130$\cdot$cos30$^\circ$. Os dois tratam a linha das cargas como se estivesse sobre o plano. A resolução declara essa idealização; o gabarito a usa sem dizer.%
\item%
\textbf{b)} Nos dois, a parte física é a mesma: a soma vale (M+m+m\_c) g sen$\theta$ e a repartição entre os apoios é estaticamente indeterminada (4 incógnitas e 3 equações independentes). A resolução vai além e propõe valores individuais sob um fechamento que ela própria declara não decorrer da estática. Esses valores não estão errados como hipótese, mas não podem ser deduzidos do enunciado. Conferi o exemplo: f\_A $\approx$ 33,44 N e f\_B $\approx$ 31,56 N, com soma 65,0 N = 130$\cdot$sen30$^\circ$. Como a resolução também afirma que só a soma é fixada, a conclusão coincide com a do gabarito.%
\end{itemize}%
\paragraph{Pontos de dúvida}%
\begin{itemize}%
\item%
\textit{(statement)} O item b pede o atrito "em cada apoio", mas, com dois apoios com atrito, o problema é estaticamente indeterminado. Não se sabe o que a banca do ITA esperava: só a soma (como o Poliedro), f\_i = N\_i$\cdot$tan$\theta$ (atrito proporcional à normal, como no fechamento da resolução) ou f\_A = f\_B. Isso pode mudar o critério de correção do item. A resolução cobre os dois primeiros casos, porque dá a soma e também a repartição sob uma hipótese declarada.%
\item%
\textit{(key)} O gabarito diz que os dois atritos apontam para cima do plano. A estática só fixa o sentido da soma; com a repartição indeterminada, um dos atritos individuais poderia apontar para baixo. Essa afirmação também é uma hipótese implícita.%
\item%
\textit{(solution)} O fechamento f\_i = N\_i$\cdot$tan$\theta$ (reação vertical em cada apoio) é arbitrário, e a própria resolução cita f\_A = f\_B como alternativa. Um corretor pode descontar se considerar que os valores individuais aparecem como resposta do item. O risco é pequeno, porque a resolução identifica claramente que eles são uma hipótese e não um resultado da estática.%
\item%
\textit{(statement)} Na figura, a barra fica elevada sobre apoios triangulares. Se a linha das cargas estiver a uma altura $\delta$ acima dos contatos, as normais mudam de $\pm$$\delta$$\cdot$g$\cdot$sen$\theta$$\cdot$(M+m+m\_c)/(L-2d); conferi que isso dá cerca de 0,41 N no exemplo com $\delta$ = 1 cm. O enunciado não fornece $\delta$, então desprezá{-}la, como fazem a resolução e o gabarito, é a leitura pretendida. Só a resolução declara essa idealização.%
\item%
\textit{(key)} O gabarito não menciona as condições de validade, que a resolução explicita: N\_A $\geq$ 0 e N\_B $\geq$ 0 (colchetes não negativos), L > 2d e $\mu$\_s $\geq$ tan$\theta$. Não afeta a resposta.%
\end{itemize}%
\section*{Questão 3}%
\href{https://poliedroresolve.sistemapoliedro.com.br/vestibulares/ita/2026/ita-2-fase-dia-3-30-10-2025-2-fase-dia-3/dissertativa-3-fisica-2-fase-dia-3-ita-2026}{Resolução do gabarito}\par\smallskip%
\begin{longtable}{@{}p{0.06\textwidth} p{0.33\textwidth} p{0.33\textwidth} p{0.1\textwidth} p{0.08\textwidth}@{}}\toprule Item & Resolução & Gabarito & Veredito & Certo \\ \midrule
a & N\_r = M{[}g cos$\theta$ + $\omega$$^2$(L/2 - a){]} = (M/2){[}2g cos$\theta$ + $\omega$$^2$(L - 2a){]}, ao longo da haste, com sentido -ê\_r (do trecho longo para o pivô e daí para o trecho curto, de comprimento a) & P\_R = Mg cos$\theta$ + M$\omega$$^2$(L - 2a)/2, ao longo da haste, apontando para a extremidade mais próxima do pivô (trecho de comprimento a) & igual & ambos \\
b & cos$\theta$ = $\omega$$^2$(2a - L)/(2g), ou seja, $\theta$ = arccos{[}$\omega$$^2$(2a - L)/(2g){]}, com existência se $\omega$$^2$|L - 2a| $\leq$ 2g & $\theta$ = arccos{[}-$\omega$$^2$(L - 2a)/(2g){]}, considerando a < L/2 & igual & ambos \\
\bottomrule\end{longtable}%
\begin{itemize}%
\item%
\textbf{a)} A expressão e o sentido são os mesmos. Refiz a conta: com d = L/2 - a, a aceleração radial de G é -$\omega$$^2$d ê\_r, e F$\cdot$ê\_r + Mg cos$\theta$ = -M$\omega$$^2$d. Logo a componente dirigida para o trecho curto é Mg cos$\theta$ + M$\omega$$^2$(L - 2a)/2. Conferi numericamente com valores aleatórios de M, L, a, $\theta$, $\omega$ e $\alpha$: diferença nula. O gabarito resolve no referencial girante, somando as forças centrífugas dos dois trechos com a massa de cada um no seu ponto médio. Isso é exato, porque a força centrífuga é linear em r: (M/2L){[}(L-a)$^2$ - a$^2${]}$\omega$$^2$ = M$\omega$$^2$(L-2a)/2. A resolução usa o teorema do centro de massa no referencial inercial. As duas abordagens são equivalentes. Na solução, a fórmula vale para qualquer a, e o sinal de N\_r indica o sentido.%
\item%
\textbf{b)} -(L - 2a) = 2a - L, então as expressões são idênticas. M se cancela nas duas. A resolução acrescenta a condição de existência |cos$\theta$| $\leq$ 1, que o gabarito omite. Também observa que, com a < L/2, $\theta$ fica em (90$^\circ$, 180$^\circ${]}: o trecho longo fica acima da horizontal, e não na configuração desenhada na figura. Isso não contradiz o gabarito.%
\end{itemize}%
\paragraph{Pontos de dúvida}%
\begin{itemize}%
\item%
\textit{(key)} O gabarito não dá a condição de existência $\omega$$^2$(L - 2a) $\leq$ 2g para que arccos tenha sentido, nem nota que, com a < L/2, o ângulo de anulação é obtuso (diferente da configuração desenhada). A resolução cobre os dois pontos.%
\item%
\textit{(key)} Na transcrição da resolução do gabarito há linhas emendadas e um '+' no lugar de '=' em '(a/L)M$\omega$$^2$(a/2) + P\_R + Mg cos$\theta$ + ...'. Parece erro de transcrição/digitação. As linhas seguintes e o resultado final estão corretos.%
\item%
\textit{(statement)} O item (b) pede o ângulo 'em função de M', mas M se cancela. Gabarito e resolução tratam isso corretamente. O enunciado também não define a convenção de sinal da 'componente radial'. Os dois a tomam como positiva no sentido do trecho curto, e a resolução declara a convenção (F$\cdot$ê\_r = -N\_r).%
\end{itemize}%
\section*{Questão 4}%
\href{https://poliedroresolve.sistemapoliedro.com.br/vestibulares/ita/2026/ita-2-fase-dia-3-30-10-2025-2-fase-dia-3/dissertativa-4-fisica-2-fase-dia-3-ita-2026}{Resolução do gabarito}\par\smallskip%
\begin{longtable}{@{}p{0.06\textwidth} p{0.33\textwidth} p{0.33\textwidth} p{0.1\textwidth} p{0.08\textwidth}@{}}\toprule Item & Resolução & Gabarito & Veredito & Certo \\ \midrule
a & y$_{2}$ = 2b ou y$_{2}$ = -4b. O CM, em (4a/5, (b+y$_{2}$)/5), é imposto sobre a elipse x$^2$/a$^2$ + y$^2$/b$^2$ = 1. & y$_{2}$ = 2b ou y$_{2}$ = -4b, com a mesma premissa de que o CM continua na elipse original. O Poliedro diz que essa premissa é falsa e sugere anular a questão. & igual & ambos \\
b & U\_min = -(GMm/5)$\cdot${[}16/$\surd$(25a$^2$ - 15b$^2$ - 24ac) + 1/$\surd$(2a$^2$ + 3b$^2$ - 2ac){]}, com c = $\surd$(a$^2$-b$^2$), estrela no foco (c,0) e y$_{2}$ = 2b. O termo mútuo -16Gm$^2$/(25$\surd$(a$^2$+49b$^2$)) aparece à parte, por ser da ordem de m/M. & E\_pg,mín = -GMm$_{2}$$\cdot${[}4/$\surd$((3a/4 - $\surd$(a$^2$-b$^2$))$^2$ + (b/4)$^2$) + 1/$\surd$((a - $\surd$(a$^2$-b$^2$))$^2$ + (2b)$^2$){]}, com y$_{2}$ = 2b e a estrela em (c,0). & equivalente & ambos \\
\bottomrule\end{longtable}%
\begin{itemize}%
\item%
\textbf{a)} Refiz a conta. Com m$_{1}$ = 4m/5 e m$_{2}$ = m/5, o CM fica em (4a/5, (b+y$_{2}$)/5). Na elipse, 16/25 + (b+y$_{2}$)$^2$/(25b$^2$) = 1, o que dá (b+y$_{2}$)$^2$ = 9b$^2$ e, portanto, y$_{2}$ = 2b ou y$_{2}$ = -4b. As duas respostas são idênticas. Ambas avisam que a premissa só vale logo após a cisão: a resolução diz que o erro é de segunda ordem em $\delta$/r e grande neste caso; o Poliedro sugere anular a questão.%
\item%
\textbf{b)} As duas expressões são a mesma escrita de outro jeito. Como m$_{2}$ = m/5, o fator -GMm$_{2}$ é igual a -GMm/5. Expandindo com c$^2$ = a$^2$ - b$^2$: d$_{1}$$^2$ = (3a/4 - c)$^2$ + b$^2$/16 = (25a$^2$ - 15b$^2$ - 24ac)/16, logo 4/d$_{1}$ = 16/$\surd$(25a$^2$ - 15b$^2$ - 24ac). Já d$_{2}$$^2$ = (a - c)$^2$ + 4b$^2$ = 2a$^2$ + 3b$^2$ - 2ac. Conferi numericamente com python3: para (a,b) = (1; 0,6), (2; 1,3) e (3; 0,4), as duas fórmulas dão -5,22404, -2,53260 e -1,34561 (em GM$\cdot$m, unidades de a), iguais até 1e-14. A escolha do mínimo é a mesma nas duas: foco (c,0) e y$_{2}$ = 2b. Também conferi a tabela de força bruta da resolução: -5,224, -5,143, -0,606 e -0,580. Estão corretos o caso circular ($\approx$ -1,10 GMm/a) e a razão |U\_mut|/|U\_*| $\approx$ 0,028$\cdot$m/M. O gabarito não inclui o termo mútuo; a resolução o deixa fora da resposta principal. Por isso as respostas principais coincidem.%
\end{itemize}%
\paragraph{Pontos de dúvida}%
\begin{itemize}%
\item%
\textit{(statement)} A questão depende de uma premissa que o enunciado não declara: depois da cisão, o CM de m$_{1}$ + m$_{2}$ continuaria na elipse original. Isso só vale logo após a separação. Com os dados, a distância P$_{1}$P$_{2}$ = $\surd$(a$^2$ + 49b$^2$)/4 é comparável a a, e aí os fragmentos sentem campos bem diferentes. O Poliedro sugere anular a questão; a resolução também aponta a falha, mas segue o modelo. Convém confirmar se a banca anulou ou alterou a questão, porque isso muda a nota dos dois itens.%
\item%
\textit{(statement)} O enunciado não diz em qual foco está a estrela nem qual dos dois valores de y$_{2}$ vale. O item b resolve isso com o 'menor energia possível': escolhe o foco (c,0) e y$_{2}$ = 2b. A resolução e o gabarito justificam a escolha da mesma forma (as duas distâncias diminuem ao mesmo tempo). No item a, dar os dois valores parece ser o esperado.%
\item%
\textit{(key)} O enunciado pede a energia potencial 'do sistema', e o gabarito usa só os termos estrela--fragmento, sem o termo mútuo -Gm$_{1}$m$_{2}$/d. Isso é razoável, porque M ? m e o termo mútuo é da ordem de m/M ($\approx$ 2,8\% de m/M no caso a=1, b=0,6). A resolução escreve esse termo à parte: -16Gm$^2$/(25$\surd$(a$^2$+49b$^2$)). Na correção, isso não deveria ser penalizado, já que ela não o soma à resposta principal.%
\item%
\textit{(key)} Na justificativa do gabarito, a frase 'y$_{2}$ = -4b deixaria m$_{2}$ mais afastada do eixo x' está certa, mas é informal. A justificativa rigorosa é que r$_{2}$$^2$ depende de y$_{2}$ apenas por y$_{2}$$^2$, e |2b| < |-4b|, como faz a resolução. Isso não muda o resultado.%
\end{itemize}%
\section*{Questão 5}%
\href{https://poliedroresolve.sistemapoliedro.com.br/vestibulares/ita/2026/ita-2-fase-dia-3-30-10-2025-2-fase-dia-3/dissertativa-5-fisica-2-fase-dia-3-ita-2026}{Resolução do gabarito}\par\smallskip%
\begin{longtable}{@{}p{0.06\textwidth} p{0.33\textwidth} p{0.33\textwidth} p{0.1\textwidth} p{0.08\textwidth}@{}}\toprule Item & Resolução & Gabarito & Veredito & Certo \\ \midrule
a & W\_ciclo = 70$\cdot$g$\cdot$$\alpha$$\cdot$L$\cdot$(M1 - M2) = 700$\cdot$$\alpha$$\cdot$L$\cdot$(M1 - M2) J (positivo se M1 > M2) & W = 70$\cdot$(M1 - M2)$\cdot$g$\cdot$$\alpha$$\cdot$L & igual & ambos \\
b & $\eta$ = (M1 - M2)$\cdot$g$\cdot$$\alpha$$\cdot$L / {[}$\rho$\_H2O$\cdot$V$\cdot$c\_H2O + 3Rm/M\_Cu + M1$\cdot$g$\cdot$$\alpha$$\cdot$L{]} & $\eta$ = (M1 - M2)$\cdot$g$\cdot$$\alpha$$\cdot$L / {[}$\rho$\_H2O$\cdot$V$\cdot$c\_H2O + 3Rm/M\_Cu + M1$\cdot$g$\cdot$$\alpha$$\cdot$L + m$\cdot$g$\cdot$$\alpha$$\cdot$L/2{]} & diferente & gabarito \\
\bottomrule\end{longtable}%
\begin{itemize}%
\item%
\textbf{a)} A expressão é a mesma. Os dois somam +M1$\cdot$g$\cdot$$\alpha$L$\Delta$T do aquecimento com -M2$\cdot$g$\cdot$$\alpha$L$\Delta$T do resfriamento, com $\Delta$T = 70 K. A resolução só substitui g = 10 m/s$^2$, o que dá 700$\alpha$L(M1-M2) em J, e explica o sinal. Também confere o resultado pela área do retângulo no plano (y, N).%
\item%
\textbf{b)} A diferença está só na parcela m$\cdot$g$\cdot$$\alpha$$\cdot$L/2 do denominador. Ela vem do aumento de energia potencial da própria viga: a base fica apoiada no piso e o centro de massa sobe $\Delta$L/2 = $\alpha$L$\Delta$T/2 no aquecimento. Pela 1ª lei, Q\_q = $\Delta$U(água + cobre) + M1$\cdot$g$\cdot$$\Delta$L + m$\cdot$g$\cdot$$\Delta$L/2, e o gabarito está correto ao incluir esse termo. A resolução trata exatamente esse efeito nas 'Verificações': diz que ele soma $\frac12$mg$\alpha$L$\Delta$T a Q\_H e Q\_F, sem alterar W, e o considera desprezível. A razão em relação ao termo de Dulong{-}Petit é gM\_Cu$\alpha$ L/(6R) $\approx$ 2,2$\times$10$^{-}$$^{7}$$\cdot$L, valor que conferi. Mesmo assim, a resposta final deixa o termo de fora. Isso é incoerente, porque a resolução mantém M1$\cdot$g$\cdot$$\alpha$$\cdot$L, que é da mesma ordem de grandeza (razão 2M1/m). No exemplo numérico, $\eta$ = 2,0117015$\times$10$^{-}$$^{8}$ sem o termo e 2,0117015$\times$10$^{-}$$^{8}$ com ele, ou seja, são numericamente iguais. Por isso a física está certa e a falha é de completude na expressão final. O resto coincide com o gabarito: W/Q\_H, água $\rho$Vc$\Delta$T, cobre 3Rm$\Delta$T/M\_Cu e trabalho M1g$\alpha$L$\Delta$T.%
\end{itemize}%
\paragraph{Pontos de dúvida}%
\begin{itemize}%
\item%
\textit{(solution)} No item b, a expressão final de $\eta$ não tem a parcela m$\cdot$g$\cdot$$\alpha$$\cdot$L/2 (energia potencial da viga, cujo centro de massa sobe $\Delta$L/2), que o gabarito inclui em Q\_q. A resolução reconhece o termo e o descarta como desprezível, mas mantém M1$\cdot$g$\cdot$$\alpha$$\cdot$L, que é da mesma ordem. Um corretor pode descontar parte da nota do item b por resposta literal incompleta.%
\item%
\textit{(key)} Os dois põem em Q\_q o trabalho mecânico (M1$\cdot$g$\cdot$$\Delta$L e m$\cdot$g$\cdot$$\Delta$L/2) além de $\Delta$U. Também seria possível interpretar como calor fornecido apenas $\Delta$U da água e do cobre. Como esses termos são numericamente desprezíveis, a escolha não muda $\eta$ na prática, mas muda a forma literal da resposta.%
\item%
\textit{(statement)} Hipóteses que nenhum dos dois declara ou que a resolução declara e o gabarito não: usar c\_V = 3R para o sólido em processo a pressão constante (c\_P $\approx$ c\_V para sólidos); desprezar a dilatação da água e o trabalho contra a atmosfera; desprezar a capacidade térmica do recipiente; supor que a troca M1$\to$M2 ocorre sem deslocamento. São todas razoáveis e iguais nas duas versões.%
\item%
\textit{(solution)} A resolução traz uma observação não pedida sobre a 2ª lei: o modelo de viga incompressível com U = U(T) e carga externa permite que $\eta$ passe do limite de Carnot. Com M1 = 10$^{9}$ kg e M2 = 0, a fórmula dá $\eta$ $\approx$ 0,80, contra 1 - 293/363 $\approx$ 0,19 (conferi as duas contas). A observação está correta e não afeta a resposta, mas foge do escopo e poderia confundir um corretor.%
\end{itemize}%
\section*{Questão 6}%
\href{https://poliedroresolve.sistemapoliedro.com.br/vestibulares/ita/2026/ita-2-fase-dia-3-30-10-2025-2-fase-dia-3/dissertativa-6-fisica-2-fase-dia-3-ita-2026}{Resolução do gabarito}\par\smallskip%
\begin{longtable}{@{}p{0.06\textwidth} p{0.33\textwidth} p{0.33\textwidth} p{0.1\textwidth} p{0.08\textwidth}@{}}\toprule Item & Resolução & Gabarito & Veredito & Certo \\ \midrule
a & L = f1 + f2 = 85 cm (0,85 m), no ajuste normal (sistema afocal) & L = f1 + f2 = 80 + 5 = 85 cm & igual & ambos \\
b & |G| = f1/f2 = 16 (G = -16 com sinal); imagem final invertida (girada de 180$^\circ$), virtual e no infinito & G = f1/f2 = 16 em módulo; com sinal, G = -16; imagem invertida & igual & ambos \\
c & Objetiva: p1' = -80 cm (virtual, direita, m1 = +2). Ocular: p2 = 165 cm, p2' = 165/32 $\approx$ 5,16 cm. Imagem final real, invertida, a $\approx$ 5,16 cm atrás da ocular ($\approx$ 90,2 cm da objetiva), m = -1/16 $\approx$ -0,0625, portanto menor que o objeto & p1' = -80 cm; p2 = 165 cm; p2' $\approx$ 5,2 cm; A $\approx$ -0,063. Imagem real, invertida, menor que o objeto, a 5,2 cm da ocular & equivalente & ambos \\
\bottomrule\end{longtable}%
\begin{itemize}%
\item%
\textbf{a)} Os dois usam o mesmo raciocínio: imagem intermediária no foco imagem da objetiva, que coincide com o foco objeto da ocular. Daí L = 80 + 5,0 = 85 cm.%
\item%
\textbf{b)} Mesmo valor (16 em módulo, -16 com sinal) e mesma orientação: invertida. A resolução também diz que a imagem é virtual e está no infinito, o que é correto e não contradiz o gabarito.%
\item%
\textbf{c)} Refiz a conta com frações. Na objetiva, p1' = 40$\cdot$80/(40-80) = -80 cm e m1 = +2. Na ocular, p2 = 85 + 80 = 165 cm, p2' = 165$\cdot$5/160 = 165/32 = 5,15625 cm e m2 = -1/32. O aumento total é m = -1/16 = -0,0625. A diferença está só no arredondamento: o gabarito usa 5,2 cm e -0,063, a resolução dá os valores exatos 165/32 cm e -1/16. As conclusões são as mesmas: imagem real, invertida e menor que o objeto.%
\end{itemize}%
\paragraph{Pontos de dúvida}%
\begin{itemize}%
\item%
\textit{(key)} Em (c), o gabarito arredonda p2' para 5,2 cm e usa esse valor arredondado no cálculo de A, obtendo -0,063. O valor exato é m = -1/16 = -0,0625. A diferença é só de arredondamento e não muda a conclusão.%
\item%
\textit{(statement)} Em (b), o sinal de G depende da convenção. O gabarito aceita G = 16 em módulo e G = -16 com sinal, e a resolução dá os dois. A orientação (invertida) é o que importa e está igual nos dois.%
\item%
\textit{(statement)} Em (c), resolução e gabarito supõem que o tubo continua com L = 85 cm (ajuste do item a), embora o enunciado não diga isso. É a leitura natural. A imagem final real fica a cerca de 5,2 cm da ocular, do lado do observador. Isso responde ao que foi pedido (posição e tipo), mesmo que essa imagem não seja observável de forma confortável.%
\end{itemize}%
\section*{Questão 7}%
\href{https://poliedroresolve.sistemapoliedro.com.br/vestibulares/ita/2026/ita-2-fase-dia-3-30-10-2025-2-fase-dia-3/dissertativa-7-fisica-2-fase-dia-3-ita-2026}{Resolução do gabarito}\par\smallskip%
\begin{longtable}{@{}p{0.06\textwidth} p{0.33\textwidth} p{0.33\textwidth} p{0.1\textwidth} p{0.08\textwidth}@{}}\toprule Item & Resolução & Gabarito & Veredito & Certo \\ \midrule
a & I = x\^{}\{N+1\} I\_0 cos\^{}\{2N\}($\varphi$/N) & I\_final = I0$\cdot$x\^{}(N+1)$\cdot$cos\^{}(2N)(?/N) & igual & ambos \\
b & r = sqrt(P/($\pi$ x\^{}\{N+1\} I\_0 cos\^{}\{2N\}($\varphi$/N))) = |cos($\varphi$/N)|\^{}\{{-}N\} sqrt(P/($\pi$ x\^{}\{N+1\} I\_0)), com cos($\varphi$/N) $\neq$ 0 & r = sqrt( P / ($\pi$$\cdot$I0$\cdot$x\^{}(N+1)$\cdot$cos\^{}(2N)(?/N)) ) & igual & ambos \\
\bottomrule\end{longtable}%
\begin{itemize}%
\item%
\textbf{a)} Mesma expressão. Os dois montam a mesma recorrência: I\_1 = x I\_0 (o eixo do 1º polarizador é paralelo à polarização do laser, então só entra a perda) e I\_k = x cos$^2$$\beta$ I\_\{k{-}1\} para k = 2,\ldots{},N+1, com $\beta$ = $\varphi$/N. Isso dá N+1 fatores de x e N fatores de cos$^2$$\beta$. Conferi em Python os exemplos numéricos da resolução e todos estão corretos: N=2, x=0,8, $\varphi$=$\pi$/2 dá 0,128 I\_0; N=2, x=0,9, $\varphi$=$\pi$/2, I\_0=1000 W/m$^2$ dá 182,25 W/m$^2$; N=10 e N=100 dão 284 e 0,024 W/m$^2$.%
\item%
\textbf{b)} Os dois partem de P = I\_final$\cdot$$\pi$r$^2$ (feixe cilíndrico e homogêneo) e chegam à mesma expressão. A segunda forma da resolução, com |cos|\^{}N fora da raiz, é equivalente. A resolução ainda trata o caso degenerado cos($\varphi$/N) = 0, que o gabarito não menciona. Isso não muda a resposta. O exemplo numérico também confere: I $\approx$ 45,17 W/m$^2$ e r $\approx$ 4,2 mm para P = 2,5 mW.%
\end{itemize}%
\paragraph{Pontos de dúvida}%
\begin{itemize}%
\item%
\textit{(statement)} O enunciado admite duas leituras do fator x: perda composta, aplicada em cada polarizador, ou um x único para o sistema inteiro. A resolução e o gabarito adotam a leitura composta, com x\^{}(N+1), que é a mais natural para "após a passagem por cada polarizador". A resolução também cita a leitura alternativa (x I\_0 cos\^{}\{2N\}($\varphi$/N)), mas deixa claro que não é a adotada. Um corretor rigoroso poderia ver nisso uma resposta "com alternativa", mas a resposta final é inequívoca.%
\item%
\textit{(solution)} Hipótese implícita, declarada na resolução e tácita no gabarito: o analisador só mede a potência e não é um polarizador adicional, ou seja, não entra nos N+1 nem aplica fator de Malus ou de perda. Isso é razoável, porque nenhum ângulo do analisador é dado.%
\item%
\textit{(key)} O gabarito não trata o caso cos($\varphi$/N) = 0, em que I = 0, P = 0 e r fica indeterminado. É só uma questão de completude e não afeta a nota.%
\end{itemize}%
\section*{Questão 8}%
\href{https://poliedroresolve.sistemapoliedro.com.br/vestibulares/ita/2026/ita-2-fase-dia-3-30-10-2025-2-fase-dia-3/dissertativa-8-fisica-2-fase-dia-3-ita-2026}{Resolução do gabarito}\par\smallskip%
\begin{longtable}{@{}p{0.06\textwidth} p{0.33\textwidth} p{0.33\textwidth} p{0.1\textwidth} p{0.08\textwidth}@{}}\toprule Item & Resolução & Gabarito & Veredito & Certo \\ \midrule
a & P\_min = 35 W (E\_bat = 36 kJ + E\_term = 6 kJ = 42 kJ em T = 1200 s) & P\_min = 35 W (36000 J + 6000 J = 42000 J em 1200 s) & igual & ambos \\
b & $\Delta$t\_max = 14L/V = 2,24$\pi$$^2$$\times$10$^{-}$$^{6}$ s $\approx$ 2,2$\times$10$^{-}$$^{5}$ s (L = $\mu$0$\cdot$n$^2$$\cdot$?$\cdot$A = 8$\pi$$^2$$\times$10$^{-}$$^{7}$ H) & $\Delta$t\_max = (56/25)$\cdot$$\pi$$^2$$\cdot$10$^{-}$$^{6}$ s $\approx$ 2,2$\cdot$10$^{-}$$^{5}$ s & igual & ambos \\
\bottomrule\end{longtable}%
\begin{itemize}%
\item%
\textbf{a)} Refiz a conta: Q = 2 Ah = 7200 C, então E\_bat = 7200$\cdot$5,0 = 36 000 J. O calor é E\_term = 100 g$\cdot$2,0 J/(g$^\circ$C)$\cdot$30 $^\circ$C = 6000 J. O total dá 42 000 J, e 42 000/1200 = 35 W. A resolução e o gabarito dão o mesmo valor com o mesmo raciocínio: o mínimo corresponde ao caso sem perdas.%
\item%
\textbf{b)} (56/25)$\pi$$^2$$\cdot$10$^{-}$$^{6}$ é exatamente 2,24$\pi$$^2$$\cdot$10$^{-}$$^{6}$ $\approx$ 2,211$\times$10$^{-}$$^{5}$ s, conferido em python. O gabarito usa N$\mu$0n$\pi$r$^2$$\cdot$14/5, que é o mesmo que $\mu$0n$^2$?A$\cdot$14/V, já que N = n? = 40, e L = 8$\pi$$^2$$\times$10$^{-}$$^{7}$ H $\approx$ 7,90 µH. A resolução também faz uma cota energética: com ciclos encadeados e E\_max = $\frac12$LI0$^2$ $\approx$ 7,74$\times$10$^{-}$$^{4}$ J por ciclo, $\Delta$t $\leq$ E\_max/P\_min. Essa cota dá o mesmo valor porque P\_min = $\frac12$$\cdot$I0$\cdot$V = 35 W. Para justificar o 'máximo', a resolução usa |M| $\leq$ L (acoplamento perfeito dá o maior $\Delta$t). O gabarito diz apenas que um $\Delta$t maior induziria menos de 5 V. As duas justificativas valem e o resultado final é igual.%
\end{itemize}%
\paragraph{Pontos de dúvida}%
\begin{itemize}%
\item%
\textit{(statement)} A prova traz $\mu$0 = 4$\pi$$\times$10$^{-}$$^{7}$ N$\cdot$A$^{-}$$^{1}$, mas a unidade certa é N$\cdot$A$^{-}$$^2$ (H/m). O gabarito copia 'N/A' sem comentar, e a resolução aponta o erro de impressão. Como o valor numérico é o mesmo, isso não muda a resposta.%
\item%
\textit{(statement)} A frase 'varia linearmente entre 0 e 14 A em um período de tempo $\Delta$t' foi lida pelos dois como uma rampa de 0 a 14 A com duração $\Delta$t. Se alguém lesse $\Delta$t como o período de uma onda triangular (sobe e desce dentro de $\Delta$t), cada subida duraria $\Delta$t/2 e o resultado dobraria ($\approx$ 4,4$\times$10$^{-}$$^{5}$ s). A leitura usual de 'período de tempo' como 'intervalo de tempo' sustenta a solução adotada.%
\item%
\textit{(solution)} A hipótese H3 da resolução diz que a energia entregue por ciclo não passa de $\frac12$LI0$^2$, a energia magnética no pico. Essa hipótese não é justificada com rigor, porque num acoplamento real a corrente do secundário também muda o fluxo. Ela só serve para a cota energética do passo 3; o resultado vem direto de Faraday, $\varepsilon$ = L$\cdot$14/$\Delta$t = V, no passo 4. Por isso não afeta a resposta, mas um corretor pode achar essa parte supérflua ou discutível.%
\item%
\textit{(key)} O gabarito e a resolução usam o modelo ideal do solenoide, B = $\mu$0nI, embora ?/d = 2 não seja um solenoide longo. A resolução estima pela fórmula de Wheeler que o L real fica cerca de 19\% abaixo ($\approx$ 6,4 µH, $\Delta$t $\approx$ 1,8$\times$10$^{-}$$^{5}$ s). É claro, porém, que a banca espera o modelo ideal.%
\end{itemize}%
\section*{Questão 9}%
\href{https://poliedroresolve.sistemapoliedro.com.br/vestibulares/ita/2026/ita-2-fase-dia-3-30-10-2025-2-fase-dia-3/dissertativa-9-fisica-2-fase-dia-3-ita-2026}{Resolução do gabarito}\par\smallskip%
\begin{longtable}{@{}p{0.06\textwidth} p{0.33\textwidth} p{0.33\textwidth} p{0.1\textwidth} p{0.08\textwidth}@{}}\toprule Item & Resolução & Gabarito & Veredito & Certo \\ \midrule
a & d = 6$\surd$2$\times$10$^{-}$$^{7}$ m $\approx$ 8,5$\times$10$^{-}$$^{7}$ m ($\approx$ 848,5 nm), de tg$\theta$$_{1}$ = 20/20 $\Rightarrow$ $\theta$$_{1}$ = 45$^\circ$ e d$\cdot$sen45$^\circ$ = 600 nm & d = 600$\surd$2 nm ($\approx$ 848,5 nm), de tg$\theta$ = 20/20 $\Rightarrow$ $\theta$ = 45$^\circ$ e d$\cdot$sen45$^\circ$ = 600 nm & equivalente & ambos \\
b & x\_V = 20$\surd$14/7 cm $\approx$ 10,7 cm, de sen$\theta$\_V = $\surd$2/3, cos$\theta$\_V = $\surd$7/3, tg$\theta$\_V = $\surd$14/7 & y' = 20$\surd$14/7 cm ($\approx$ 10,7 cm), de sen$\theta$' = $\surd$2/3 & igual & ambos \\
\bottomrule\end{longtable}%
\begin{itemize}%
\item%
\textbf{a)} Mesmo valor, só muda a unidade: 6$\surd$2$\times$10$^{-}$$^{7}$ m = 600$\surd$2 nm. A resolução também dá o valor em nm (848,5 nm). O método é o mesmo nos dois: tg$\theta$ = x/H para achar o ângulo e sen $\theta$ na equação da grade. Conferi numericamente: 600/sen45$^\circ$ = 848,53 nm.%
\item%
\textbf{b)} Mesmo resultado. O gabarito resolve y'/$\surd$(D$^2$+y'$^2$) = $\surd$2/3, que dá 7y'$^2$ = 2D$^2$. A resolução calcula tg$\theta$ a partir de sen$\theta$. Os dois caminhos são equivalentes. Conferi numericamente: sen $\theta$' = 400/848,53 = 0,4714 = $\surd$2/3 e 20$\cdot$tg(arcsen 0,4714) = 10,690 cm = 20$\surd$14/7. A resolução também aponta que a aproximação de ângulos pequenos (que daria 13,3 cm) não vale aqui, porque $\theta$ = 45$^\circ$.%
\end{itemize}%
\paragraph{Pontos de dúvida}%
\begin{itemize}%
\item%
\textit{(statement)} O enunciado fala em "projeção do primeiro máximo\ldots{} a 20 cm do ponto central". Tanto a resolução quanto o gabarito entendem isso como distância horizontal no solo, o que dá tg$\theta$ = 1. Se alguém tomasse 20 cm como o comprimento do raio difratado, teria sen$\theta$ = 1, o que não faz sentido fisicamente. Por isso a leitura adotada é a única razoável, e a ambiguidade não muda a resposta.%
\item%
\textit{(solution)} A resolução declara hipóteses que o gabarito deixa implícitas: grade de transmissão, meio com n $\approx$ 1, luz solar tratada como feixe colimado (a largura angular do Sol, cerca de 0,5$^\circ$, é desprezada) e ordem m = 1 do mesmo lado. Nenhuma delas muda o resultado.%
\end{itemize}%
\section*{Questão 10}%
\href{https://poliedroresolve.sistemapoliedro.com.br/vestibulares/ita/2026/ita-2-fase-dia-3-30-10-2025-2-fase-dia-3/dissertativa-10-fisica-2-fase-dia-3-ita-2026}{Resolução do gabarito}\par\smallskip%
\begin{longtable}{@{}p{0.06\textwidth} p{0.33\textwidth} p{0.33\textwidth} p{0.1\textwidth} p{0.08\textwidth}@{}}\toprule Item & Resolução & Gabarito & Veredito & Certo \\ \midrule
a & $\Delta$t = t\_B = L(c+v1)/{[}v1(c+v2){]} = 387/32 anos = 12,09375 anos $\approx$ 12,1 anos (tempo coordenado no referencial da Terra, contado a partir de t0) & t\_Terra = 10,75 + 1,34375 = 12,09375 anos $\approx$ 12,1 anos & igual & ambos \\
b & $\Delta$$\tau$ = t\_B/$\gamma$2 = 0,8 $\cdot$ 12,09375 = 387/40 anos = 9,675 anos $\approx$ 9,7 anos (tempo próprio da nave 2); confirmado pela transformação de Lorentz do evento de emissão (t'\_A = 6,9875; x'\_A = 2,6875; t'\_B = 9,675) & t\_nave2 = t\_Terra/$\gamma$2 = 12,09375/(5/4) = 9,675 anos & igual & ambos \\
\bottomrule\end{longtable}%
\begin{itemize}%
\item%
\textbf{a)} Refeito com frações exatas: t\_A = 8,6/0,8 = 43/4 = 10,75 anos; a nave 2 está em 6,45 al, a 2,15 al atrás; com aproximação c + 0,6c = 1,6c, $\Delta$t = 2,15/1,6 = 1,34375 anos; logo t\_B = 387/32 = 12,09375 anos. A resolução chega ao mesmo valor por uma fórmula fechada, igual à do gabarito. A verificação x\_B = 7,25625 al < L também confere.%
\item%
\textbf{b)} $\gamma$2 = 1/$\surd$(1-0,36) = 5/4 e $\tau$ = 12,09375 $\cdot$ 4/5 = 9,675 anos. Os eventos O e B ficam na linha de universo da nave 2, então o intervalo é tempo próprio. Conferi também pela transformação de Lorentz do evento A (10,75; 8,6) para S': t'\_A = 6,9875 anos e x'\_A = 2,6875 al; somando o trânsito da luz, dá t'\_B = 9,675 anos.%
\end{itemize}%
\paragraph{Pontos de dúvida}%
\begin{itemize}%
\item%
\textit{(statement)} 'Tempo transcorrido na Terra' poderia ser lido como o instante em que a Terra recebe a luz da recepção, o que daria 12,09375 + 7,25625 = 19,35 anos. A resolução e o gabarito usam a coordenada de tempo do evento no referencial da Terra, que é a leitura usual em relatividade. A resolução comenta essa alternativa e a descarta.%
\item%
\textit{(statement)} Hipóteses implícitas que os dois adotam: a distância de 8,6 al é medida no referencial da Terra e contada a partir da fronteira do Sistema Solar; Sirius está em repouso em relação à Terra; as 'trajetórias paralelas' são tratadas como colineares, sem separação transversal; e a emissão do sinal é imediata. Nenhuma delas muda a resposta.%
\end{itemize}%
\end{document}